\tocsection{\textbf{9/8/26 Lecture 4: }Estimation and Sufficiency}

Recall the setting. We observe data $X \sim P$ for some $P \in \cal{P} = \{P_{\theta} : \theta \in \Theta\}$. We want to guess estimand $g(\theta)$ (which is typically the median, mean, or some coefficient) by applying an estimator $\delta(\cdot)$ to $X$.\\

A good guess is something that minimizes a function $L(\theta, d)$, which is called the loss. This measures the disutility of a guess. Typically, we take $L(\theta, d) = (g(\theta) - d)^2$. Finally, to tell how good a estimator is given any sample, we use the risk function
\begin{align*}
    R(\theta; \delta) &= \bb{E}_{\theta}[L(\theta, \delta(X))]\\
    &= \int_{X} L(\theta, \delta(X))dP_{\theta}(X)
\end{align*}
for example, an example of the risk function is
\[ \tt{MSE}(\delta) = \bb{E}_{\theta}[(g(\theta) - \delta(X))^2] \]

For instance, consider $X \sim \tt{Binom}(n, \theta)$, which has probability mass function $p_{\theta}(x) = {n \choose x}\theta^{x}(1 - \theta)^{n - x}$. We want to estimate $g(\theta) = \theta$. Consider the following estimators
\[ \delta_0(X) = \frac{X}{n},\quad \delta_1(X) = \frac{X + 1}{n + 2},\quad \delta_2(X) = \frac{X + 2}{n + 4},\quad \text{ and } \delta_{3}(X) = \frac{X + 1}{n} \]
Notice that although the first one is unbiased, the others have smaller variance, so their MSEs are comparable. Here is a graph of their MSEs for a given $\theta$. 

\fig{0.6}{Lecture4Fig1}{}

We see that between $\theta_0, \theta_1, \theta_2$, there is always one $\theta$ that makes one estimator better than another. However, we see that $\theta_3$ is strictly worse than $\theta_0$, which means there is no reason to use it over $\theta_0$. We will soon formalize what this means, but first: Why don't we just consider an estimator that is the best among all estimators?\\

Well, the reason is that it doesn't exist. If there exists an estimator that is always better than every estimator everywhere, then for a given $\theta$, we can take an estimator that simply gives the constant estimate $\theta$. Then, it has no $\tt{MSE}$, which implies that our ``best everywhere'' estimator must have no $\tt{MSE}$ everywhere, which simply cannot happen. 

\begin{definition}[\textbf{Definition 4.1:} Inadmissable Estimators]
    We say an estimator $\delta$ is \emph{inadmissable} if there exists an estimator $\delta^{*}$ such that
    \begin{enumerate}
        \item $R(\theta; \delta^{*}) \leq R(\theta; \delta)$ for all $\theta \in \Theta$.
        \item $R(\theta; \delta^{*}) < R(\theta; \delta)$ for some $\theta \in \Theta$. 
    \end{enumerate}
    We then say that $\delta^{*}$ strictly dominates $\delta$. 
\end{definition}

Thus, in our case above, $\delta_{3}$ is inadmissable.\\

Now, given multiple estimators that are inadmissable, how do we remove the ambiguity. The way to do this is to reduce the risk function $R(\theta; \delta)$ to a scalar (not dependent on $\theta$). There are multiple ways to do this.\\

\textbf{Method 1: Average Case Risk:} This involes measuring the average case risk, which means we are going to average over $\theta$. Thus, we need to assume that $\theta \sim \pi$ and minimize $\int_{\theta} R(\theta;\delta) d\pi(\theta)$. If $\pi$ is a probability measure, then we are trying to find $\bb{E}_{\pi}[R(\theta; \delta)]$. For instance, if $\pi$ is uniform on $[0, 1]$ in the case above, then we summarize the risk functions by integrating over the MSE over $[0, 1]$.\\

Notice that this isn't inherently Bayes. We aren't assuming that $\theta$ takes a distribution, but we are using this distribution as a way to average $\theta$. However, if we are doing Bayesian statistics, then this would be the way to pick the best estimator.\\

\textbf{Method 2: Worse Case Risk:} With this method, we are trying to minimize, over $\theta$, the $\sup_{\theta \in \Theta} R(\theta; \delta)$. From out example above $\delta_2$ minimizes worse case risk among the other estimators. \\

\textbf{Method 3: Restrict Estimators:} One thing we can do is to just rule out all unbiased estimators and try Method 1 or Method 2 on the resulting list. This has merit in some cases, but we will see sometimes it's useful to use biased estimators.

\subsection*{Sufficiency}
\sectionspace

Recall the three models of coin flips: \\

\textbf{Model 3:} $X_{i, t} \overset{\tt{ind}}{\sim} \tt{Bern}(\theta_{i, t})$, for $i \leq 48$ and $t \leq n_{i}$\\

\textbf{Model 2:} $X_{i} \overset{\tt{ind}}{\sim} \tt{Binom}(\theta_{i})$ for $i \leq 48$. Notice that $X_{i} \sum_{t = 1}^{n_{i}} X_{i, t}$ where $X_{i, t} \overset{\tt{ind}}{\sim} \tt{Bern}(\theta_{i})$.\\

\textbf{Model 1:} $X \overset{\tt{ind}}{\sim} \tt{Binom}(n^{+}, \theta)$. Notice that $X = \sum_{i} X_{i} = \sum_{i, t} X_{i, t}$. 

Thus, in some sense
\[ \text{Model 1} \subseteq \text{Model 2} \subseteq \tt{Model 3} \]
as given that we know $X_{i, t}$ but Model 2 is our assumption, we can construct $X_{i}$, and given that we know $X_{i}$ but Model 1 is our assumption, we can construct $X$.

\begin{definition}[\textbf{Definition 4.2:} Sufficient Statistics]
    A statistics $T(X)$ is \emph{sufficient} for a model $\cal{P}$ if the conditional distribution of $X | T(X)$ is the same for all $P \in \cal{P}$. In other words, no matter what $\cal{P}$ is, knowing that $T(X) = c$ can tell you what the distribution of $X$ is. 
\end{definition}

For example, if $X_1, \ldots, X_n$ are i.i.d. $\tt{Bern}(\theta)$, then $T(X) = \sum_{i} X_i$ is sufficient, because knowing that $T(X) = c$ can give you the marginal distribution of $X = (X_1, \ldots, X_n)$.\\

Now, notice that $X_i$ is sufficient for $\tt{Binom}(n_i, \theta_i)$, it is not sufficient for $\tt{Bern}(\theta_{i, t})$, and is not minimal sufficient for $\tt{Binom}(n, \theta)$ (as it gives too much information).

\subsection*{Factorization Theorem}
\begin{theorem}[\textbf{Theorem 4.3:} Factorization Theorem]
    Let $\cal{P} = \{P_{\theta} \col \theta \in \Theta\}$ be a model with densities $p_{\theta}(X)$ with respect to a common dominating measure $\mu$ on $X$. Then, $T(X)$ is sufficient if, and only if, there exists a $g_{\theta}(t)$ and $h(x) \geq 0$ such that
    \[ p_{\theta}(x) \overset{\mu \tt{ a.e.}}{=} g_{\theta}(T(x)) \cdot h(x) \]
\end{theorem}
\begin{proof}{Proof:}
    We will prove it for discrete $X$, as the proof for the continuous case is similar but just requires measure theoretic notions. First, assume wlog that $\mu = \#$ (as every measure is absolutely continuous with respect to $\#$). \\

    $(\Leftarrow)$ We see that
    \begin{align*}
        \bb{P}_{\theta}(X = x \given T = t) &= \frac{\bb{P}_{\theta}(X = x, T(x) = t)}{\bb{P}_{\theta}(T(x) = t)}\\
        &= \frac{\bb{P}_{\theta}(X = x, T(x) = t)}{\sum_{z \col Tz = t} \bb{P}_{\theta}(X = z)}\\
        &= \frac{g_{\theta}(t)h(x) \cdot 1\{T(x) = t\}}{\sum_{T(z) = t} g_{\theta}(t)h(z)}\\
        &= \frac{h(z)1\{T(x) = t\}}{\sum_{T(z) = t}h(z)}
    \end{align*}
    Thus, we see that this doesn't depend on $\theta$, so is the same for all $P \in \cal{P}$.\\

    $(\Rightarrow)$ Assume that $T$ is sufficient. Then, let $g_{\theta}(t) = \bb{P}_{\theta}(T(X) = t)$ and $h(x) = \bb{P}_{\theta}(X = x \given T(x) = t)$. Then, as
    \begin{align*}
        g_{\theta}(T(x))h(x) &= \bb{P}_{\theta}(T(X) = T(x) \text{ and } X = x)\\
        &= \bb{P}_{\theta}(X = x)\\
        &= p_{\theta}(x)
    \end{align*}
    which was to be shown.

    $\qed$
\end{proof}
\textbf{Note:} We can absorb $h(x)$ into $\mu$ by letting $\nu$ such that $\frac{d\nu}{d\mu} = h$. Thus, we would then have $p_{\theta}(X) = g_{\theta}(T(X))$, which is what we kind of mean by ``$T$ gives the distribution of $X$''.\\

For example, if $X_1, \ldots, X_n \overset{\tt{iid}}{\sim} \cal{N}(\theta, 1)$, then as the density function is $\phi(x) = \frac{1}{\sqrt{2\pi}}e^{-\frac{1}{2}(x - \theta)^2}$, we see that
\begin{align*}
    p_{\theta}(x) &= (2\pi)^{-n/2} \prod_{i = 1}^{n} e^{-1/2(x_i - \theta)^2}\\
    &= e^{\theta\sum_{i = 1}^{n} x_i - n\theta^2/2} \cdot \frac{e^{-\sum x_i^2/2}}{(2\pi)^{n/2}}
\end{align*}
which is the desired factorization showing that $\sum_{i = 1}^{n} x_i$ is sufficient.\\

Another example is that if $X_1, \ldots, X_n \overset{\tt{iid}}{\sim} \tt{Pois}(\theta)$, then as $\phi(x) = \frac{{\theta}^{x}e^{-\theta}}{x!}$ for $x = 0, 1, \ldots$, then as
\[ p_{\theta}(x) = \prod_{i = 1}^{n} \frac{\theta^{x_i}e^{-\theta}}{x_i!} = \theta^{\sum x_i}e^{-n\theta} \cdot \frac{1}{\prod_{i = 1}^{n} x_i!} \]
we see that $\sum_{i = 1}^{x} x_i$ is sufficient.\\

One last interesting example is that, if $X_i \overset{\tt{iid}}{\sim} U[\theta, \theta + 1]$, then as $\theta(x) = 1\{\theta \leq x \leq \theta + 1\}$, we see that
\[ p_{\theta}(x) = \prod_{i = 1}^{n} 1\{\theta \leq x_i \leq \theta + 1\} = 1\{\theta \leq \min_{i} X_i \leq \max_{i} X_i \leq \theta + 1\} \]
we can conclude that $(\min_{i} X_i, \max_{i} X_i)$ is sufficient.
